\documentclass[a4paper,12pt]{scrartcl}

% Layout
\usepackage[a4paper, left=2.5cm, right=2cm, top=2cm, bottom=3.2cm]{geometry} % kleinere Ränder
\usepackage{parskip}

% Umlaute in der Datei erlauben, auf Deutsch umstellen
\usepackage{fontspec}
\usepackage{lmodern}
\usepackage[english, ngerman]{babel}

% Mathesymbole und Ähnliches
\usepackage{amsmath}
\usepackage{amsthm}
\usepackage{mathtools}
\usepackage{amssymb}
\usepackage{microtype}
\usepackage{stmaryrd}
\usepackage{MnSymbol}
\usepackage{bussproofs}		% Sequenzenkalkül
\usepackage{xcolor}

% Grafiken und PDFs einfügen
\usepackage{graphicx}
\usepackage{pdfpages}
\usepackage{float}

% PDF-Tools
\usepackage[hidelinks, unicode]{hyperref}

% Abbildungen
\usepackage{tikz}
\usetikzlibrary{arrows,calc,positioning}
\tikzset{ 
   state/.style={circle,fill,inner sep=2pt,label=below:#1\strut},
   arrow/.style={->, shorten <=3pt, shorten >=3pt, >=stealth',thick},
   knoten/.style={draw, circle, minimum size=8mm}} 
   
% Infoboxen 
\definecolor{orange}{RGB}{230,159,0}
\definecolor{darkgray}{RGB}{70,90,70}

\usepackage{tcolorbox}


% Reelle, Natürliche, Ganze, Rationale Zahlen
\newcommand{\R}{\ensuremath{\mathbb{R}}}
\newcommand{\N}{\ensuremath{\mathbb{N}}}
\newcommand{\Z}{\ensuremath{\mathbb{Z}}}
\newcommand{\Q}{\ensuremath{\mathbb{Q}}}

% Fraktur für Strukturen
\newcommand{\A}{\ensuremath{\mathfrak A}}
\newcommand{\B}{\ensuremath{\mathfrak B}}
\newcommand{\I}{\ensuremath{\mathfrak I}}

% Körper
\newcommand{\K}{\ensuremath{\mathcal{K}}}
%\newcommand{\G}{\ensuremath{\mathcal{G}}}

% Makros für logische Operatoren
\newcommand{\xor}{\ensuremath{\oplus}} % exklusives oder
\newcommand{\impl}{\ensuremath{\rightarrow}} % logische Implikation
\newcommand{\Lor}{\ensuremath{\bigvee}}
\newcommand{\Land}{\ensuremath{\bigwedge}}

\newcommand{\mal}{\cdot}

% Meistens ist \varphi schöner als \phi, genauso bei \theta
\renewcommand{\phi}{\varphi}
\renewcommand{\theta}{\vartheta}
\renewcommand{\epsilon}{\varepsilon}

\newcommand*{\todo}{\textbf{\textcolor{red}{TODO}}}		% big, red, apparent TODO

% SEQUENZENKALKÜL
\usepackage{bussproofs}
\newcommand{\yields}{\ensuremath{\mathrel{\Rightarrow}}}

% Aufzählungen anpassen (alternativ: \arabic, \alph)
\renewcommand{\labelenumi}{(\roman{enumi})}

% BITTE NICHT ÄNDERN: interne Kommandos für die Informationen (Blattnummer, Gruppe, ...)
\newcounter{sheet}
\newcommand{\printsheet}{\arabic{sheet}}

\newcommand{\printgroup}{?}
\newcommand{\group}[1]{%
\renewcommand{\printgroup}{#1}%
}

\newcommand{\printmembers}{}
\newcommand{\printmember}[3]{{#2} {#3} & {#1} \\}

% Für Vektoren
\newcommand{\myvec}[1]{\ensuremath{\begin{pmatrix}#1\end{pmatrix}}}
\def\spvec#1{\left(\vcenter{\halign{\hfil$##$\hfil\cr \spvecA#1;;}}\right)}
\def\spvecA#1;{\if;#1;\else #1\cr \expandafter \spvecA \fi}

\newtheorem{proposition}{Proposition}[sheet]

\newcommand{\pdfmembers}{}
\newcommand{\pdfmember}[3]{{#1} {#3}, {#2}; }

\newcommand{\member}[3]{%
\expandafter\renewcommand\expandafter\printmembers\expandafter{\printmembers\printmember{#1}{#2}{#3}}%
\expandafter\renewcommand\expandafter\pdfmembers\expandafter{\pdfmembers\pdfmember{#1}{#2}{#3}}%
}

\AtBeginDocument{\hypersetup{
    pdftitle = {Übungsblatt \printsheet},
    pdfauthor = {Abgabegruppe \printgroup: \pdfmembers},
    pdfsubject = {Online Algorithms}
}}

% <-------- HIER alle Informationen eintragen ========================
\setcounter{sheet}{3} % Nummer des Blatts

% alle Gruppenmitglieder in der Form \member{Matrikelnummer}{Vorname}{Nachname}
\member{5481114}{Danje}{Petersen}

\begin{document}

% Kopfzeile
{\raggedright
\begin{tabular}{l}
    Online Algorithms \\
    SS 2024 \\
\end{tabular}}
\hfill
{\Large Sheet 0\printsheet}
\hfill
\begin{tabular}{l l}
\today{} \\
\printmembers
\end{tabular}
\hrule


% <-------- HIER beginnt die Lösung ========================

\section*{Exercise \printsheet.1: Online Bin Packing - Next Fit}
For the first $n+1$ bins we first give an item of size $\alpha$ and then give items such that there is only space $\alpha-\varepsilon$ left in the bin and Next Fit needs to continue with the next bin. After $n$ bins we give $n+1$ requests for items of size $\alpha-\varepsilon$. Next Fit will fill one bin to exactly $1$ before it fills each of the following bins with $\left\lfloor \frac{1}{\alpha-\epsilon} \right\rfloor$ items, therefore using  $\frac{n}{\left\lfloor\frac{1}{\alpha-\epsilon}\right\rfloor}$ bins. \\
Overall, Next Fit needs to use at least $n+n\mal (\alpha-\epsilon)$ bins, while OPT only needs $n$ which leads to: 
\[\frac{c_{NF}}{c_{OPT}} \geq \frac{n+ n \mal (\alpha-\epsilon)}{n} = 1+ \alpha -\epsilon \] 

For $\alpha \leq \frac{1}{2}$ the offset $\epsilon$ can be very small which makes $\alpha - \epsilon$ almost as big as $\alpha$ and therefore
\[\frac{c_{NF}}{c_{OPT}} \geq 1 + \alpha \] %\frac{1}{\alpha}\]

For $\alpha > \frac{1}{2}$ the offset is such that $\alpha - \epsilon = 1-\alpha$. In this case: 
\[\frac{c_{NF}}{c_{OPT}} \geq 1+ 1-\alpha\]

\section*{Exercise \printsheet.2}
\subsection*{(a)}

\subsection*{(b)}


\section*{Exercise \printsheet.3: Online Graph Coloring}
Let A be any deterministic online algorithm for Graph Coloring. For any $k \in \N$ we construct a forest $T_k$ that forces A to use $k$ colors.\\
$T_1$ is just a single vertex that obviously needs to be colored in one color.
For $T_2$ we construct the Graph $T_1$ and connect it to a new vertex, forcing A to use a new color. \\

In the same manner we construct a forest $T_k$ for $k\in \N$ recursively. First we construct $T_i$ for all $i \in \{1,\dots, k-1\}$. Since in each $T_i$ there must be at least one more color than in $T_{i-1}$, A already used $k-1$ colors. Now we introduce a new vertex that is connected to each $T_i$ via the vertex that uses a color that is not used in $T_{i-1}$. Since that new vertex has neighbours with all the $k-1$ colors already used by A, the algorithm needs to use a new color. \\

Since each $T_i$ is a forest and these are connected through a single vertex, there is only one unique path from each $T_i$ to another, making $T_k$ a forest too.







\end{document}
